\part{Project Management}
% LaTeX Canvas: Chapter Outline
\chapter{Project Management}
\label{ch:project-management}

\section*{Purpose}
Most student projects fail for non-technical reasons: files scattered across desktops, unclear goals, missing data provenance, undocumented assumptions, and work tracked in private messages rather than a shared system. This chapter introduces lightweight project management practices that make small data-science projects stable, auditable, and collaborative.

\section*{Learning objectives}
By the end of this chapter, you should be able to:
\begin{enumerate}
\item Define a project goal, success criteria, and constraints in a short project brief.
\item Set up a reproducible project structure with clear locations for data, code, notebooks, outputs, and docs.
\item Apply data management hygiene: provenance, naming, immutability of raw data, and data dictionaries.
\item Write minimal project documentation that enables another person (or future you) to reproduce results.
\item Use issue tracking to plan work, document decisions, and coordinate collaboration.
\item Maintain a simple cadence for planning, execution, review, and delivery.
\end{enumerate}

\section*{Running theme: make progress visible and work repeatable}
A well-managed project makes its assumptions and status legible: what the data are, what changed, what remains, and how to rerun.

\section{Mental model: a project is a system of artifacts}
\subsection{Artifacts you must manage}
\begin{itemize}
\item \textbf{Data:} raw inputs, processed outputs, intermediate files.
\item \textbf{Code:} scripts, reusable modules, notebooks.
\item \textbf{Environment:} dependencies and runtime assumptions.
\item \textbf{Documentation:} purpose, decisions, usage, and interpretation.
\item \textbf{Work tracking:} tasks, bugs, questions, and priorities.
\end{itemize}

\subsection{The lifecycle: plan $\rightarrow$ build $\rightarrow$ verify $\rightarrow$ deliver}
\begin{itemize}
\item \textbf{Plan:} decide what success looks like.
\item \textbf{Build:} implement data pipelines/analysis.
\item \textbf{Verify:} check quality, reproduce outputs.
\item \textbf{Deliver:} package results and communicate limitations.
\end{itemize}

\section{Project planning for novices (lightweight, not bureaucratic)}
\subsection{The one-page project brief}
\begin{itemize}
\item Problem statement (one paragraph).
\item Audience and use case.
\item Deliverables (tables, plots, memo, dashboard, model).
\item Success criteria (what ``done'' means).
\item Constraints (time, compute, access, privacy).
\item Risks and unknowns (data availability, definitions, biases).
\end{itemize}

\subsection{Decompose work into milestones}
\begin{itemize}
\item Milestone 1: data acquisition + intake note.
\item Milestone 2: cleaning + data dictionary.
\item Milestone 3: analysis + validation.
\item Milestone 4: outputs + narrative + reproducibility check.
\end{itemize}

\subsection{Define decision points}
\begin{itemize}
\item What choices might change results? (filters, joins, missingness handling)
\item Where will you record those decisions? (issues, changelog, notebook narrative)
\end{itemize}

\section{Reproducible project structure}
\subsection{Design principles}
\begin{itemize}
\item One project folder = one project.
\item Raw data are read-only.
\item Outputs are reproducible from code.
\item Paths are relative to the project root.
\item Clear separation: data vs code vs docs vs outputs.
\end{itemize}

\subsection{A student-friendly directory template}
\begin{verbatim}
project/
README.md
environment.yml   (or requirements.txt)
data/
raw/
processed/
external/
notebooks/
src/
scripts/
reports/
figures/
docs/
tests/            (optional)
\end{verbatim}

\subsection{Naming conventions and consistency}
\begin{itemize}
\item Use lowercase and hyphens/underscores (choose one).
\item Date stamps for snapshots (e.g., \texttt{2026-01-08}).
\item Avoid spaces and ambiguous names (\texttt{final\_final2.csv}).
\end{itemize}

\subsection{What goes where (and what does not)}
\begin{itemize}
\item \texttt{data/raw}: original sources; never edited.
\item \texttt{data/processed}: generated outputs from cleaning.
\item \texttt{notebooks}: narrative drivers; keep tidy.
\item \texttt{src}: reusable functions and modules.
\item \texttt{scripts}: command-line entry points.
\item \texttt{reports/figures}: final artifacts for sharing.
\item Avoid storing secrets or credentials anywhere in the repo.
\end{itemize}

\section{Data management hygiene}
\subsection{Provenance: where did this data come from?}
\begin{itemize}
\item Source URL or file origin.
\item Retrieval date/time.
\item License/terms of use.
\item Any access restrictions or privacy considerations.
\end{itemize}

\subsection{Raw data immutability}
\begin{itemize}
\item Treat raw data as evidence: do not edit.
\item If you must fix an obvious error, do it in a processing step that is documented.
\end{itemize}

\subsection{Data dictionary and codebook}
\begin{itemize}
\item One table with: variable name, type, meaning, units, allowable values, missingness codes.
\item Record transformations: derived fields, recodes, joins.
\end{itemize}

\subsection{Versioning data}
\begin{itemize}
\item When datasets change, create dated snapshots or version tags.
\item Record checksums or file sizes for integrity (intro concept).
\item Keep ``what changed'' notes (new rows, new columns, schema shifts).
\end{itemize}

\subsection{Sensitive data hygiene (baseline)}
\begin{itemize}
\item Identify whether data contain identifiers.
\item Minimize access: store sensitive data outside repos when required.
\item Never commit secrets, API keys, or tokens.
\item Redact or aggregate before sharing.
\end{itemize}

\section{Documentation that makes a project runnable}
Chapter~\ref{ch:documentation} covers how to find, read, and write technical documentation in depth; this section focuses on the documentation artifacts specific to a project.
\subsection{README as the single source of truth}
\begin{itemize}
\item What the project does (one paragraph).
\item What data are required and where they go.
\item How to set up the environment.
\item One ``copy/paste'' command to reproduce core outputs.
\item How to interpret outputs (where to find the key tables/figures).
\end{itemize}

\subsection{Decision log}
\begin{itemize}
\item Record consequential choices: filtering rules, outlier handling, join keys.
\item Store decisions in issues or a dedicated \texttt{DECISIONS.md}.
\end{itemize}

\subsection{Docstrings and inline comments (just enough)}
\begin{itemize}
\item Explain intent, not obvious syntax.
\item Document inputs, outputs, assumptions.
\end{itemize}

\subsection{Reproducibility notes}
\begin{itemize}
\item Environment file is required.
\item Note OS-specific dependencies.
\item Include a small ``smoke test'' script or notebook cell.
\end{itemize}

\section{Issue tracking fundamentals (for students)}
\subsection{Why issues are not just for bugs}
\begin{itemize}
\item Tasks, questions, data problems, and design decisions.
\item A shared memory and audit trail.
\end{itemize}

\subsection{Anatomy of a good issue}
\begin{itemize}
\item Clear title: action + object.
\item Description: context, objective, definition of done.
\item Evidence: links, screenshots, error messages.
\item Labels: type (bug/task/question), priority, area (data/code/docs).
\end{itemize}

\subsection{Issue workflow: triage $\rightarrow$ do $\rightarrow$ review $\rightarrow$ close}
\begin{itemize}
\item Triage: clarify, label, assign, and break down.
\item Do: create a branch or working copy tied to the issue (concept).
\item Review: verify outputs, update docs.
\item Close: summarize what changed and link artifacts.
\end{itemize}

\subsection{Milestones and project boards (optional)}
\begin{itemize}
\item Milestones map to deliverables.
\item Boards visualize status (To do / Doing / Done).
\item Keep it lightweight: avoid over-customization early.
\end{itemize}

\section{Quality gates: make ``done'' meaningful}
\subsection{The reproducibility check}
\begin{itemize}
\item Fresh environment or clean restart.
\item Run the pipeline end-to-end.
\item Confirm outputs match expectations.
\end{itemize}

\subsection{Data checks}
\begin{itemize}
\item Row/column counts.
\item Missingness summary.
\item Value ranges and plausibility.
\item Join coverage and key uniqueness.
\end{itemize}

\subsection{Output checks}
\begin{itemize}
\item Figures labeled, units clear, captions present.
\item Tables have consistent formatting and definitions.
\item Narrative explains limitations and uncertainty.
\end{itemize}

\section{Common failure modes and how to prevent them}
\subsection{Folder chaos and lost files}
\begin{itemize}
\item Prevention: stable project root + fixed structure.
\item Recovery: search by extension/date; reconstruct from outputs.
\end{itemize}

\subsection{Silent data drift}
\begin{itemize}
\item Prevention: snapshot raw data; record retrieval dates and checksums.
\item Detection: schema checks and summary stats comparisons.
\end{itemize}

\subsection{Undocumented assumptions}
\begin{itemize}
\item Prevention: decision log + notebook narrative.
\item Practice: ``if this changed, results would change'' gets documented.
\end{itemize}

\subsection{Work tracked in private channels}
\begin{itemize}
\item Prevention: issues for tasks/questions; link decisions and artifacts.
\end{itemize}

\section{Worked examples (outline)}
\subsection{Example 1: Start a new course project in 20 minutes}
\begin{itemize}
\item Create the folder structure.
\item Add environment file.
\item Write a README with ``how to run''.
\end{itemize}

\subsection{Example 2: Intake a dataset with provenance and a data dictionary}
\begin{itemize}
\item Place raw file.
\item Write provenance notes.
\item Draft codebook and a first-pass quality report.
\end{itemize}

\subsection{Example 3: Use issues to manage a cleaning pipeline}
\begin{itemize}
\item Create issues for missingness, type conversions, duplicates, joins.
\item Close each issue with a summary and link to outputs.
\end{itemize}

\subsection{Example 4: Final reproducibility check before submission}
\begin{itemize}
\item Recreate env.
\item Run end-to-end.
\item Confirm outputs and update README.
\end{itemize}

\section{Templates}
\subsection{Template A: One-page project brief}
\begin{verbatim}
Title:
Problem:
Audience:
Deliverables:
Success criteria:
Constraints:
Risks/unknowns:
Milestones:
\end{verbatim}

\subsection{Template B: README skeleton}
\begin{verbatim}

# Project name

## Purpose

## Data

* Source:
* Retrieved:
* License/notes:
* Location: data/raw/

## Setup

* Create environment:
* Activate environment:

## Run

* Command(s) to reproduce key outputs:

## Outputs

* reports/
* figures/

## Notes

* Decisions and limitations
  \end{verbatim}

\subsection{Template C: Issue template (student version)}
\begin{verbatim}
Title:
Type: bug/task/question
Context:
What I tried:
Evidence (errors, screenshots, links):
Definition of done:
\end{verbatim}

\section{Exercises}
\begin{enumerate}
\item Create a new project folder using the template and write a README that someone else could follow.
\item Intake a dataset: place it in \texttt{data/raw}, write provenance notes, and draft a data dictionary.
\item Create five issues that decompose the project into milestones and tasks; label and prioritize them.
\item Implement one task and close its issue with a short ``what changed'' summary.
\item Perform a reproducibility check by recreating your environment and rerunning the pipeline.
\end{enumerate}

\section{One-page checklist}
\begin{itemize}
\item I have a clear project goal and definition of done.
\item My project folder has a reproducible structure.
\item Raw data are immutable and provenance is recorded.
\item I maintain a data dictionary and transformation notes.
\item My README explains setup, run, and outputs.
\item Work is tracked in issues with labels and clear closure notes.
\item I can reproduce results from a clean environment.
\end{itemize}

% \appendix
\section{Quick reference: ``minimum viable'' project operations}
\begin{itemize}
\item Create structure.
\item Write README.
\item Record environment.
\item Intake data with provenance.
\item Track work in issues.
\item Reproduce end-to-end before delivery.
\end{itemize}
